\begin{center}
  {\heiti\zihao{4} 摘要}
\end{center}
\vspace{0.5em}

\noindent
大模型 Token 已经成为智能服务定价、路由、缓存与治理的基本操作单位。但是，Token 发票并不等同于福利账户：相同数量的 Token 既可能只生成一段普通回答，也可能减少一次行政办事错误、节省专业人员劳动时间，或引入新的风险。本文围绕中国“人工智能+”国家战略，提出 TESSA-PSA（Token Economy Statistical Satellite Accounting for Public-Service Allocation）统计建模框架，用于度量和配置 Token 驱动的公共智能服务能力。与宏观叙事式的 Token 经济核算不同，本文将建模核心压缩为一个可落地的问题：给定任务库、公开 API 价格、模型能力记录、人工节时标注与 Token 预算，估计福利调整后的边际 Token 生产率，并在公共服务与企业支持任务之间配置稀缺智能服务容量。方法上，本文结合难度校准的 Token 生产模型、随机化或双重稳健的反事实价值估计、就业与工资等易获得宏观总量的校准，以及带有风险和公平约束的保守预算优化。数据设计强调竞赛可执行性，所需数据可来自公开价格页、中文开放基准、政府 FAQ、人工标注和官方工资就业统计。本文的创新不在于提出新的提示词优化技巧，而在于构造一套可复现的“测度—估计—配置—看板”协议，把国家战略叙事转化为可估计指标、可视化证据与可操作的资源配置规则。

\vspace{1em}
\noindent{\heiti 关键词：}Token 经济；人工智能+；公共智能服务；因果估值；资源配置
